% Derived integration fragment, 2026-10-08. No documentclass.
% Numerical results in the companion JSON are floating diagnostics only.
\subsection{A parameterized one-step double sieve at a fixed total level}
\label{sec:generalized-one-step-double-sieve}

This extends the directly proved one-step construction in
\texttt{same-new-level-double-sieve-route.tex}; it does not transfer an
old $H,h$ table. The published combinatorial input is
\href{https://www.impan.pl/shop/en/publication/transaction/download/product/82961}
{Wu 2004, Lemma 4.1, (4.1)--(4.5), and (5.1)--(5.10)}.
The signed non-switched distribution input is the existing prefix
interface, using
\href{https://arxiv.org/html/2109.02851v2}{Lichtman 2021, (5.14)--(5.18)}
and
\href{https://arxiv.org/html/2309.08522v1}
{Lichtman 2023, Propositions 5.2 and 6.1}, with the already proved
selected-subproduct repair. The switched input is the separately
derived positive-count theorem with $K=681/100$. None of these uses
identifies that positive-count conclusion with the signed switched
distribution hypothesis in L21 (5.22).

Write $b=p_1\cdots p_r$, with decreasing specified exponents
$t_1>\cdots>t_r>0$, and set
\[
 T=\sum t_i,\quad Q=N^\nu,\quad R=Q/b=N^{L+o(1)},\quad
 L=\nu-T>0,\quad \alpha=\frac{1-T}{L},\qquad
 A(s)=\frac{sF(s)}{2e^\gamma},\quad a(s)=\frac{sf(s)}{2e^\gamma}.
\]
Here $a(s)=0$ for $s\le2$. Every formula below is the limit of
fixed positive exponent boxes, fixed level buffers, and the actual
residual weights at that one total level $Q$.

\paragraph{The precise non-switched family.}
We use either of the following two verified interfaces. In the enhanced
case $\nu>1/2$, require the original full-prefix square conditions
\begin{equation}\label{eq:generalized-one-step-squares}
 t_1+\cdots+t_{j-1}+2t_j<\nu\quad(1\le j\le r),
\end{equation}
and a factorization branch that remains valid after appending primes
below the smallest specified prime. The branches used below are the
largest-prime branch $v(t_1)$; a selected-subproduct branch involving
only the original specified slots, with all its remaining-square
gates; and a $W$ branch whose first three slots or upper bound on its
third slot persist. All require strict cap and gate margins.
The full level is below $153/256<3/5$.

If the current physical cutoff is $z=N^{L/s}$, we require
\begin{equation}\label{eq:generalized-one-step-cutoff}
 s\ge2,\qquad L/s\le t_r.
\end{equation}
For the added prime $q=N^{Lu}$ in Wu's negative term,
$1/s_0\le u\le1/s$, the new square condition is
$T+2Lu\le T+L=\nu$. The previous square conditions are unchanged.
Every subsequent residual rough square condition follows by
multiplying the residual condition by $bq$; its cutoff is
$y=R^{1/s_0}\le q$. Thus these are checks of the complete concatenated
vector, not just its bare specified product. For a selected branch,
the remaining-square tests for old unselected slots are unchanged;
before any new slot all original slots have already occurred, so
each new selected-branch test is the complete square test above.
The original signed atom coefficients, factorials, and small-prime
weight are retained by the box interface. This proves the signed
upper estimate for $S(\mathcal A_{N,b},y)$ and signed lower estimate
for $S(\mathcal A_{N,bq},y)$ at the same $Q$.
Where its lower argument is at most two, count nonnegativity supplies
the multiplier zero, without a distribution claim for that lower term.

In the classical case take $\nu<1/2$ with fixed margin. The complete
moduli of both residual sieve sums have support at most $N^\nu$.
Every specified prime is at least a fixed positive power $N^\epsilon$.
Consequently the number of ways any modulus can contain these fixed
labelled prefix slots is bounded by a constant depending only on
$r,\epsilon$. The residual modified weights have an absolute
coefficient bound depending only on their fixed construction
precision. Summing them at each complete modulus therefore leaves a
bounded coefficient, and the absolute Bombieri--Vinogradov remainder
applies, with residue $N$ on coprime moduli. The finitely many sieve
atoms are counted with their original signs; their number depends on
the fixed precision, not on $N$. This case does not require
\eqref{eq:generalized-one-step-squares}. In particular it covers the
classical part of $U_3$ even when $2t>1/2$.

\begin{proposition}[Explicit one-step upper defect]
\label{prop:generalized-one-step-defect}
On a compact positive prefix box satisfying one of the preceding
interfaces, suppose $2\le s<s_0\le5$, $\alpha\ge2$, and
\eqref{eq:generalized-one-step-cutoff}. Define
\begin{align}
 I(s,s_0,\alpha)
 &=\int_{1/s_0<u_1<u_2<u_3<1/s}
 \omega\!\left(\frac{\alpha-u_1-u_2-u_3}{u_2}\right)
                  \frac{du_1\,du_2\,du_3}{u_1u_2^2u_3},\label{eq:generalized-one-step-I}\\
 J(s,s_0,\alpha)
 &=A(s)-A(s_0)+\frac12\int_{1/s_0}^{1/s}
       \frac{a(s_0(1-u))}{u(1-u)}\,du
                   -\frac K8 I(s,s_0,\alpha).
                   \label{eq:generalized-one-step-J}
\end{align}
The ordinary normalized upper multiplier $A(s)$ can be replaced by
$A(s)-\max(0,J(s,s_0,\alpha))$. Any finite choice of $s_0$ on a common
positive prefix partition can be optimized box by box. In the limit
one can use the supremum over a fixed compact set of choices. This is
an averaged prefix estimate, not a bound for an individual prefix.
\end{proposition}

\begin{proof}
For a box indicator $\sigma(b)$, or a finite nonnegative constant
combination of such labelled indicators, put
\[
 \Phi_s=\sum_b\sigma(b)S(\mathcal A_{N,b},R^{1/s}),\qquad
 \Theta_Q=4\operatorname{li}(N)
       \sum_b\frac{\sigma(b)\mathcal C(bN)}{\varphi(b)\log(Q/b)}.
\]
Wu's pointwise identity and its prime-square error give
\[
 2\Phi_s\le2\sum_b\sigma(b)S(\mathcal A_{N,b},y)
 -\sum_b\sigma(b)\sum_{y\le q<z}S(\mathcal A_{N,bq},y)
 +\sum_b\sigma(b)\sum_{y\le q_1<q_2<q_3<z}
            S(\mathcal A_{N,bq_1q_2q_3},q_2)+o(\Theta_Q).
\]
The two non-switched inputs just proved yield, after division by
$\Theta_Q$, respectively $2A(s_0)$ and the negative of the integral
in \eqref{eq:generalized-one-step-J}. The lower argument is precisely
$s_0(1-u)$, not an argument computed at a different total level.

The positive switched multiset has the exact labelled form
\[
 m=bq_1q_2q_3n\le N,\quad y\le q_1<q_2<q_3<z,\quad P^-(n)\ge q_2.
\]
Canonically order the prime factors of $n$, retain all prefix labels,
and split into its finitely many possible factor counts. Every factor
is at least $N^{\epsilon_1}$ for some fixed $\epsilon_1>0$: the old
prefix is on a positive compact box, the new primes have exponent at
least $L/s_0$, and the quotient primes have the still larger threshold
$q_2$. All order, box, roughness, and total-product constraints are
finite monomial inequalities. Repeated factors have a power-saving
mass bound. The unit quotient has mass at most
$N^{T+3L/s+o(1)}=o(N)$ by a fixed power, since
$\alpha-3/s\ge1/2$. Thus it is not an omitted boundary atom.
The ordinary rough-number mass calculation gives the actual integral
\eqref{eq:generalized-one-step-I}, with the actual box value of
$\alpha$, rather than a supremum over $\alpha$. Its smallest Buchstab
argument is at least $\alpha s-3\ge1$. Any argument-one corner is
removed by a thin strip before taking limits.

For its exact unsifted mass $X$, the positive switching theorem gives
$(K+o(1))\mathcal C(N)X/\log N$. The partial summation in Wu (5.8),
with the same prefix measure, gives
\[
 X=(I(s,s_0,\alpha)+o(1))N
              \sum_b\frac{\sigma(b)}{\varphi(b)\log(Q/b)}.
\]
Using $\mathcal C(N)\le\mathcal C(bN)$ shows that the switched count
divided by $\Theta_Q$ is at most $KI/4+o(1)$. Dividing the weighted
inequality by two proves \eqref{eq:generalized-one-step-J}; the cost
is $KI/8$. No residual-level factor is inserted into $K$.
The required extension of the original fixed-cutoff statement of the
positive switching theorem is proved immediately below.
\end{proof}

\paragraph{Why the same $K$ permits every fixed positive factor cutoff.}
Here the original cutoff $100/1149$ must be reduced; it cannot simply
be reused. Fix any $\epsilon_1>0$ before $N\to\infty$. In the proof of
\texttt{uniform-prime-products-upper.tex}, replace its $a$ by
$\epsilon_1$ and its factor bound by $k_0=\lfloor1/\epsilon_1\rfloor$.
The modulus weights and the density integral in
\texttt{extended-prefix-upper.tex} are unchanged. The following checks
give the needed extension of that proof.

Apply the $k_{\rm HB}=3$ Heath--Brown identity to each original prime
slot at its own scale. There are at most $6k_0$ variables. Choose
$12\zeta<\epsilon_1/2$, in addition to all previous modulus margins,
and then choose the smaller distribution and local-scale precisions.
Before Mellin separation discard each dyadic rectangle incompatible
with its original prime-slot scale. The product of each group's at
most six lower endpoints is then at least $N^{\epsilon_1}/O(1)$,
including on the separated rectangular extension. If no individual
variable is in $[x^{2\zeta},x^{1/3+\zeta/2}]$, every original group
must contain a long variable: otherwise its product is at most
$x^{12\zeta+o(1)}$. There are at most two long variables globally.
Thus the no-middle-variable case is impossible for three or more
original groups. For one or two groups, the same one-long and
two-long estimates apply; the small complementary products have
the original bounds $x^{10\zeta+o(1)}$ and $x^{20\zeta+o(1)}$.
In the middle-variable case the individual coefficient is $1$ or
$\mu$, up to a logarithm and Mellin twist; its rough Siegel--Walfisz
input and the divisor-bounded complement remain the original ones.
Increasing the requested logarithmic saving absorbs the fixed larger
factor count, number of rectangles, and Mellin losses.

The exact coprime-mean replacement costs
$O(X_{\rm out}N^{-\epsilon_1}\log^C N)$ by the same union bound over
prime factors of a product. Repeated primes cost
$N^{1-\epsilon_1+o(1)}$, and prime powers introduced by the identities
cost $N^{1-\epsilon_1/2+o(1)}$; the divisor-moment treatment of the
progression errors preserves these savings. Products below
$N^{1-\kappa}$ are discarded as before. Smooth outer approximations
retain the lower cutoff $N^{\epsilon_1}/2$, the labels, and the exact
canonical quotient multiplicities. At fixed positive exponent boxes
their excess mass tends to zero. This proves the same positive-count
coefficient $K$ for the present finite product families. It asserts
neither uniformity as $\epsilon_1\downarrow0$ nor distribution for an
arbitrary arithmetic reweighting. In particular no shrinking Wu
short-box class is being substituted for these fixed boxes.

\paragraph{A single Buchstab transfer to the two-prime lower terms.}
Define a nonnegative chosen correction $j(s,\alpha)$ from the
proposition on $2\le s\le3$, and put it equal to zero outside that
interval. On a stable prefix family with cutoff $N^a$ and residual
argument $s=L/a$, use this transfer only for $2\le s<4$:
\begin{equation}\label{eq:generalized-one-step-h}
 h_1(s,\alpha)=\int_{\max(2,s-1)}^3
 j\left(w,\alpha+\frac{\alpha-1}{w}\right)\frac{dw}{w}.
\end{equation}
The normalized lower multiplier $a(s)$ can be replaced by
$a(s)+h_1(s,\alpha)$. To prove this, use the exact least-prime
Buchstab identity from cutoff $R^{1/4}$ to $N^a=R^{1/s}$.
The starting lower multiplier is $a(4)$. For a subtracted prime
$q=N^{L/(w+1)}$, the child residual argument at cutoff $q$ is $w$,
and its size ratio is
\[
 \alpha_{bq}=\frac{1-T-L/(w+1)}{L-L/(w+1)}
            =\alpha+\frac{\alpha-1}{w}.
\]
The normalization changes by $dw/w$. The ordinary part is exactly
$a(4)-\int_{s-1}^3 A(w)\,dw/w=a(s)$, with the usual lowered-cutoff
upper convention $A(w)=1$ for $w<2$. The corrected child part is
\eqref{eq:generalized-one-step-h}. Its primes satisfy
$q\le N^a\le p_r$, and its square condition is $2\log_Nq\le L$.
The second added prime used inside the child's one-step inequality
is at most $q$, with its own residual square condition. Thus both
generations retain the original stable cap. This is one finite tree;
no abstract $h,H$ iteration is being assumed. We choose zero extra
lower correction for $s<2$ and $s\ge4$.

\paragraph{The actual eleven-term domains.}
Take the open parameter range in \texttt{free-weight-thresholds.tex}:
$1/12<a<h<d<e=57/224$, $w_*=25/224<h<\min(25/128,2a)$,
$1/3<h+d<1/2$, and $4a+d,4h+d<1$.
The following are additive gains in the complete Chen coefficient,
already including its factor $1/4$.

For $L_2$ use $T=t$, $\nu=v(t)=(1+t)/2$, and cutoff $p=N^t$.
Here $2t<\nu$ and the largest-prime branch persists; $\alpha=2$.
The gain is
\[
 \Delta c_{L_2}=\int_{w_*}^h
 \frac{j((v(t)-t)/t,2)}{t(v(t)-t)}\,dt.
\]
Only $t\ge1/7$ contributes when $j$ is supported on $[2,3]$.

For $U_3,U_4$ the cutoff is lowered, when necessary, to exponent
$\min\{a,(\nu-t)/2\}$. Their $j$ argument is
$\max\{2,(\nu-t)/a\}$ and their ratio is
$(1-t)/(\nu-t)$. Use $\nu=v(t)$ for $a<t<e$, and for the
$U_3$ tail $e<t<1/3$ use the classical endpoint $\nu=1/2$ approached
from below. On the enhanced branch $2t<v(t)$; the physical cutoff
does not exceed $a\le t$. The gain integrand in either case is
\[
 \frac{j(\max\{2,(\nu-t)/a\},(1-t)/(\nu-t))}{t(\nu-t)}.
\]
Integrate over $a<t<1/3$ for $U_3$ and $a<t<d$ for $U_4$.

For $L_5,L_6$ use exactly the current two-prefix cap
\[
 \nu=V_{2,a}(t,u)=
 \begin{cases}
 \max\{v(t),v(t+u),153/256-3a/8\},&t\le25/128,\\
 v(t),&t>25/128.
 \end{cases}
\]
On the part $L=\nu-t-u\ge2a$, the first square test is
$2t<v(t)\le\nu$ and the second follows from
$L\ge2a>h\ge u$, namely $t+2u<\nu$. The cutoff in the
starting lower term is $R^{1/4}\le N^a$, and every added prime in
the exact Buchstab difference is also at most $N^a$. The original
two slots therefore remain the largest. The $v(t+u)$ branch keeps
its leading selected pair, and the $W$ branch keeps third exponent
at most $a$. It still has the gate supplied by $v(t)$, exactly as
in \texttt{double-prefix-level.tex}. Empty or short rough atoms use
the same classical fallback as that proof. Thus the gain is
\[
 \Delta c_{L_i}=\iint_{\Omega_i}
 \frac{h_1(L/a,(1-t-u)/L)}{tuL}\,du\,dt\quad(i=5,6),
\]
where $\Omega_5=\{a<u<t<h\}$ and
$\Omega_6=\{a<u<h<t<d\}$; the integrand is chosen zero outside
$2\le L/a<4$. Each cap is at most $(1+t+u)/2$, so $\alpha\ge2$.
For its constant branch this follows from
$153/256-3a/8\le1/2+a$, valid already for $a\ge25/352<1/12$.

For the two four-prefix upper terms retain all eligibility gates and
the cap $\nu=L_4(t)$ of \texttt{four-prefix-safe-minimum.tex}.
Choose cutoff exponent $\min\{w,(\nu-T)/2\}$ for the new upper
option. All four initial square tests hold by that construction;
all original selected slots, selected-branch gates, and the first
three slots for the $B(z)$ branch persist after either one-step
extension. Here $T\ge4a>1/3$, so $\nu\le153/256<(1+T)/2$.
Use the smaller corrected kernel inside the existing minimum:
\[
 \Delta c_{U_i}=\frac14\int_{\Omega_i}
 \left[\min\!\left\{\frac Kz\omega\left(\frac{1-T}{z}\right),
                         \frac4{\nu-T}\right\}
 -\min\!\left\{\frac Kz\omega\left(\frac{1-T}{z}\right),
 \frac{4(1-j(\max\{2,(\nu-T)/w\},(1-T)/(\nu-T)))}{\nu-T}\right\}
 \right]d\mu(t),\quad i=10,11.
\]
The integrand is zero off the eligible set. The original ordinary
normalized multiplier is one here because $a>1/12$, as already
proved for that minimum. This replacement avoids adding overlapping
gains from its switching and ordinary options.

Every positive contribution just used has a uniform fixed minimum
prime exponent on compact boxes. For the searched choices $s_0\le5$,
the four-prefix residual satisfies $L>w\ge a$; the corrected
two-prefix children have residual at least $4a/3$; and single-prefix
residuals are at least $1/6$ in the classical tail. For the displayed
numerical parameter tuple one can choose $\epsilon_1=a/6$ after
sufficiently small buffers. This is below the old fixed cutoff but
is exactly covered by the preceding Heath--Brown adaptation.

All boundaries are first trimmed. Take $N\to\infty$ with boxes,
branch choices, and buffers fixed; then remove the sieve and smooth
approximation errors, shrink the boxes and level buffers within the
strict domains, and finally exhaust those domains. A common finite
positive partition licenses every minimum and every choice of
$s_0$. The equality sets have zero integral mass. A numerical grid
interpolation alone is not a lower certificate for these limiting
integrals.

\paragraph{Floating feasibility, not a new certified constant.}
The companion \texttt{probe\_generalized\_one\_step.py} searches a finite
$s_0$ grid in $[3,5]$, tabulates $2\le s\le3$, $2\le\alpha\le8$,
and chooses zero correction outside that $s$ range. Its JSON records
coarse and fine Gauss quadratures, interpolation meshes, and the
separate gains. The baseline parameter tuple is
\[
 (a,h,d)=(.08455645453148589,.15049959287189102,
                           .2352266768449549).
\]
The current ordinary plus four-prefix-minimum diagnostic is
$c\approx1.87414321$. The coarse mesh gives an extra
$\Delta c\approx .01598683$, and the fine mesh gives $.01600361$.
In the fine calculation the contributions are
\[
\begin{array}{c|rrrrrrr}
 &L_2&U_3&U_4&L_5&L_6&U_{10}&U_{11}\\\hline
 \Delta c&.000136371&.013460255&0&.000094742&.002272783&
 .000039405&.000000051
\end{array}
\]
The resulting diagnostic is $1.89014682$, while $K\log(1/.749)$
is $1.96820097$; the remaining gap is about $.07805415$.
The zero-margin transfer diagnostic is $b\approx .75763420$.
The floating evaluation clamps a roundoff value just below two back
to the analytically proved bound $\alpha\ge2$, avoiding an erroneous
out-of-grid zero at the exact $\alpha=2$ boundary.
Thus this one-step family supplies genuine
analytic options but the sampled gain alone is far too small for
the requested frontier. The search is not an upper bound on all
possible one-step choices or on further finite recursions. Neither
the baseline nor these floating integrals certify a new rounded
counting coefficient; the previous certified constants are unchanged.
